\documentclass[ngerman,a4paper,12pt,headsepline]{scrreprt}

%----------------- PDF CONFIG ----------------- %
\pdfinfo{    
     /Title (PDF-Titel) 
     /Subject   (PDF-Thema)    
     /Author  (Vorname Nachname) 
     /Keywords   (Stichwort1,Stichwort2)      
} 

\title{Testframework}
\author{Can Kedik}
\date{\today}



%----------------- PAKETE INKLUDIEREN ----------------- %

\usepackage{geometry} % Packet für Seitenrandabständex und Einstellung für Seitenränder
\usepackage[ngerman]{babel} % deutsche Silbentrennung

\usepackage{booktabs} %entzerrt die Tabellenzeilen und bietet verschieden dicke Unterteilungslinien
\usepackage{longtable}% Tabellen können sich nicht über mehrere Seiten 
\usepackage{graphicx} % kann LaTeX Grafiken einbinden

\usepackage{mathabx} 
\usepackage[applemac]{inputenc} % Umlaute unter Mac werden automatisch gesetzt
\usepackage[T1]{fontenc} % Zeichenencoding
\usepackage{lmodern} % typographische Qualität 
\frenchspacing % Schaltet den zusätzlichen Zwischenraum ab
\usepackage{fix-cm}
\usepackage{varioref}
\usepackage{hyperref} % verwandelt alle Kapitelüberschriften, Verweise aufs Literaturverzeichnis und andere Querverweise in PDF-Hyperlinks
\usepackage{cleveref}
\usepackage{color}
\usepackage{url}
\usepackage[babel, german=quotes]{csquotes}
\usepackage{array}
\usepackage{siunitx}
\usepackage[backend=biber,style=alphabetic]{biblatex}
\bibliography{literatur}


\usepackage[nottoc]{tocbibind}



% für Listings
\usepackage{listings}
\lstset{numbers=left, numberstyle=\tiny, numbersep=5pt, stepnumber=4, keywordstyle=\color{black}\bfseries\itshape, stringstyle=\ttfamily,showstringspaces=false,basicstyle=\footnotesize,captionpos=b}
\lstset{language=java}



%----------------- FARBEN DEFINIEREN ----------------- %
\definecolor{gray}{gray}{0.95} % Listingsbackground

%----------------- LAYOUT SETZEN ----------------- %
\geometry{left=2cm, right=2cm, top=2.5cm, bottom=2cm}
\linespread {1.25}\selectfont %1.25 da er von Haus aus 1.2 ist und 1,25 * 1,2 = 1,5 isch




%-------##-------##-------##------- ANFANG INHALT -------##-------##-------##-------%
\begin{document}



\pagenumbering{roman} % Seitennummer

%----------------- DECKBLATT -----------------%
 %----------------- KONFIGURATION ----------------- %
\pagestyle{empty} % enthalten keinerlei Kopf oder Fuß 

%----------------- HDA FBI Logo ----------------- %
\begin{figure}[t]
	\centering
	\includegraphics[width=0.6\textwidth]{Abb/logo_fbi}
\end{figure}

%----------------- INHALT ----------------- %

\begin{center}
\Large Hochschule Darmstadt \\
\normalsize \textsc{- Fachbereich Informatik -} \\

% Whitespace
\vspace{105 pt}

\Huge Name der Arbeit \\ 
\normalsize
\vspace{20 pt}

Abschlussarbeit zur Erlangung des akademischen Grades \\ 
Bachelor of Science (B.Sc.) 

\vspace{75 pt}


vorgelegt von \\
\vspace{5 pt}
Can Kedik \\
731620
\vspace{115 pt}

\begin{tabular}[h]{p{4cm}l l}
	Referent: & Prof. Dr. Ronald C. Moore\\
	Korreferent: & Prof. Dr. Bettina Harriehausen-M"uhlbauer
\end{tabular}


\end{center}

 
%----------------- ABSTRACT -----------------%
 %----------------- KONFIGURATION ----------------- %
\pagestyle{empty} % enthalten keinerlei Kopf oder Fuß


\chapter*{Abstract} % (fold)
\label{cha:abtract}
	Lorem ipsum dolor sit amet, consetetur sadipscing elitr, sed diam nonumy eirmod tempor invidunt ut labore et dolore magna aliquyam erat, sed diam voluptua. At vero eos et accusam et justo duo dolores et ea rebum. Stet clita kasd gubergren, no sea takimata sanctus est Lorem ipsum dolor sit amet. Lorem ipsum dolor sit amet, consetetur sadipscing elitr, sed diam nonumy eirmod tempor invidunt ut labore et dolore magna aliquyam erat, sed diam voluptua. At vero eos et accusam et justo duo dolores et ea rebum. Stet clita kasd gubergren, no sea takimata sanctus est Lorem ipsum dolor sit amet.

\begin{quote}
Dies ist ein Zitat.
\end{quote}

verstand, scheinen nun doch vorueber zu Dies ist der Text sein. \\
siehe: http://janeden.net/die-praeambel

 
 
%----------------- VERZEICHNISSE -----------------%
\tableofcontents % Inhaltverzeichnis

% ----- Abbildungen ----- %
%\addcontentsline{toc}{section}{Abbildungsverzeichnis} % falls in Inhalsverzeichnis
\listoffigures

% ----- Tabellen----- %
% \addcontentsline{toc}{section}{Tabellenverzeichnis}  % falls in Inhalsverzeichnis
% \fancyhead[L]{Abbildungsverzeichnis / Abkürzungsverzeichnis} %Kopfzeile links
\listoftables

% ----- Listings ----- %
% Listingverzeichnis soll im Inhaltsverzeichnis auftauchen
% \addcontentsline{toc}{section}{Listingverzeichnis}
% \fancyhead[L]{Abbildungs- / Tabellen- / Listingverzeichnis} %Kopfzeile links
\renewcommand{\lstlistlistingname}{Listingverzeichnis}
\lstlistoflistings

\pagestyle{headings} % zurueck setzen von roemische seitenanzahl


%----------------- KAPITEL : EINLEITUNG  ----------------- %
\pagenumbering{arabic}
\chapter{Einf"uhrung}
\label{cha:einfuehrung}

Lorem ipsum dolor sit amet, consetetur sadipscing elitr, sed diam nonumy eirmod tempor invidunt ut labore et dolore magna aliquyam erat, sed diam voluptua. At vero eos et accusam et justo duo dolores et ea rebum. Stet clita kasd gubergren, no sea takimata sanctus est Lorem ipsum dolor sit amet. Lorem ipsum dolor sit amet, consetetur sadipscing elitr, sed diam nonumy eirmod tempor invidunt ut labore et dolore magna aliquyam erat, sed diam voluptua. At vero eos et accusam et justo duo dolores et ea rebum. Stet clita kasd gubergren, no sea takimata sanctus est Lorem ipsum dolor sit amet.




%----------------- KAPITEL : FRAMEWORK  ----------------- %
\chapter{Das Framework}
\label{cha:framework}
\section{Verwendete Hardware} \label{sec:verwendeteHardware}
Die in \cref{sec:testbeschreibung} beschriebenen Versuche wurden mit folgender Hardware durchgef"uhrt. %
\begin{itemize}
\item Zwei Switches
\item Vier Raspberry Pis der ersten Generation
\item Ein Raspberry Pi der zweiten Generation
\item Mehrere Ethernet Kabel
\end{itemize}
Diese Ger"ate wurden in einem eigenst"andigend Netzwerk zusammengeschaltet. So sind an jedem Switch zwei Pis %
der ersten Generation angeschlossen w"ahrend an einem der Switches der Pi der zweiten Generation angeschlossen %
ist (siehe \cref{fig:AufbauVomNetzwerk}). Welche Software auf den Pis verwendet wurde, wird in \cref{sec:aufbauSoftware} erkl"art. 

\section{Aufbau der Software} \label{sec:aufbauSoftware}
Das Testframework besteht aus drei verschiedenen Teilen. %
\begin{itemize}
\item Zabbix %
\item Eigentwicklung auf dem Server %
\item Eigentwicklung auf dem Agent %
\end{itemize}
Die Eigentwicklungen sind alle in Bash programmiert, w"ahrend Zabbix eine bereits fertige Open Source L"osung %
ist. In den folgenden Abschnitten werde ich einen Einblick in diese Teile geben. %

\subsection{Zabbix}
Zabbix ist ein Open Source Netzwerk Monitoring System. Die erste Version wurde von Alexei Vladishev entwickelt, welches %
seit 2005 von der Firma Zabbix SIA weiterentwickelt wird \autocite{zabbix:Web}. % 
Ein weiterer bekannter Vertreter der Netzwerk Monitor Systeme ist Nagios, welches wie Zabbix unter der GPL Lizenz vertrieben wird \autocite{wiki:Nagios}. % 
Beide Systeme basieren auf einer Client-Server Architektur. Im weiteren wird jedoch nur Zabbix betrachtet. %
Zabbix besteht aus zwei Komponenten. %

\begin{description}
\item[Zabbix Server]Der Server hat eine auf PHP basierende Weboberfl"ache "uber die es f"ur den Benutzer m"oglich ist die Agents zu %
konfigurieren. So k"onnen manuell die Templates erstellt werden, die den Zabbix Agents mitteilen, welche Informationen dem %
Server zu "ubermitteln sind. %
"Uber die API Schnittstelle ist es m"oglich Programme f"ur den Server zu schreiben und diese auszuf"uhren. %
Der Zabbix Server ist ein sich selbst "uberwachender Agent. Dieser ist jedoch nicht beteiligt %
im allgemeinen Prozess des Frameworks. %
Trotzdem wird in den folgenden Kapiteln der Zabbix Server nicht aussen vor gelassen und in die %
den Tests betrachtet werden, dabei werden haupts"achlich die von Zabbix generierten Werte zur Perfomance und zum %
Traffic auf dem Ethernet Port genommen. Der Wert Perfomance betrachtet zwei Dinge, einmal die sogenannte Queue, welche %
angibt wie viele von den Agents "ubermittelten Werten darauf warten vom Zabbix Server verarbeitet zu werden und die %
die verarbeiteten Werte pro Sekunde. Diese Werte lassen einen Schluss auf die Leistung des Servers schlie"sen. %
Sollte die Queue einen bestimmten schwellwert erreichen wird auf dem Zabbix Dashboard eine Warnung ausgegeben das der %
Server mit der Verarbeitung nicht hinterherkommt. Dies kann soweit gehen das die gesammelten Daten auf dem Server nicht %
aktuell oder sogar fehlerhaft sind. Jedoch ist dieses Szenario in diesem Framework unwahrscheinlich, da die gr"o"se des %
Frameworks "uberschaubbar ist. Der Traffic auf dem Ethernet Port Eth0 wird betrachtet, da alle von den Agents gesammelten Informationen %
"uber diesen an den Server gelangen. Dabei wird der eingehende sowie auch der ausgehende Traffic betrachtet. 

\item[Zabbix Agent]Die Clients, die im Netzwerk "uberwacht werden sollen sind die sogenannten Agents, die an den Server %
die Informationen weiterleiten, die vom Server gefordert werden. In den sp"ateren Kapiteln wird der Hauptfokus auf der %
I/O--Last der Festplatte und wie beim Server dem ein/- und ausgehenden Traffic auf dem Ethernet Port Eth0 liegen. %
\end{description} 

\subsection{Eigenentwicklung}
Die selbstentwickelte Software wird in zwei Kategorien unterteilt; ein Teil der Software l"auft %
auf den Agents, diese haben den Zweck Netzwerk Traffic zu erzeugen. Dadurch entsteht auf dem Netzwerk und %
auf den Zabbix Agents eine Nutzlast, die vom Zabbix Server gesammelt werden kann. Der zweite Teil der Software %
die auf dem Server l"auft, die Software auf dem Server unterst"utzt den Entwicklungsprozess auf den Agents. %
\begin{enumerate}
\item Zabbix Server 
\begin{description}
\item[Update Script:]Dieses Script wird von der Software Entwicklungsmaschine ausgef"uhrt. Es aktualisiert den Code auf den Endger"aten, %
die die Programme Hintergrundrauschen, Pinger, Synchronize und Startrauschen aktualisieren. Dabei wird mittels Secure Copy %
der Quellcode auf das Endger"at gespielt.

\item[Get logs:]Die Skripte Pinger und Hintergrundrauschen erstellen jeweils auf den Endger"aten logfiles. %
Da es jedoch ein sehr hoher Verwaltungsaufwand w"are auf den Endger"aten die Logfiles weiterzuverwerten, %
werden mit dem Skript Get Logs die auf den Agents gelagerten Logfiles auf dem Rechner der dieses Skript %
startet gesammelt. Somit hat man die von den Endger"aten gesammelten Daten auf einem Rechner und kann %
dessen Weiterverarbeitung betreiben.      

\item[Calculate Average:]Berechnet den Durschnitt der Dauer, die ein Paket braucht, um erfolgreich versendet %
zu werden.
 
\end{description}

\item Zabbix Agent

\begin{description}
\item[Hintergrundrauschen:]Diese Eigentwicklung stellt mit Zabbix die Kernkomponente des Frameworks dar. Wie %
der Rest der selbstentwickelten Software wurde sie komplett in Bash programmiert, da das Framework ausschlie"slich %
in einer Linux Umgebung entwickelt, getestet und verwendet wird. Hintergrundrauschen schickt "uber Secure Copy, %
Pakete von einem Agent zum anderen. So wird eine %
Last auf dem Netzwerk erzeugt, die Mithilfe des Zabbix Servers gemessen werden kann. Ausserdem speichert Hintergrundrauschen %
die Dauer, bis ein Paket erfolgreich bei seinem zuf"allig ausgew"ahlten Empf"anger angekommen ist. %

\item[Startrauschen:]wird automatisch zum Start eines der Endger"ate ausgef"uhrt. Es dient als eine Zeitschaltuhr %
und erm"oglicht ein zeitversetztes Starten des Scripts Hintergrundrauschen, wodurch man den sequentiellen %
Anstieg an Last im Netzwerk beobachten kann.

\item[Synchronize:]Raspberry Pis besitzen keine eigene Batterie wie es handels"ubliche Rechner haben, %
deshalb ist nach jedem Neustart die Uhrzeit der Pis unzuverl"assig. Dieses Programm synchronisiert %
die Uhrzeiten der Agents mit der Zeit des Servers. %

\item[Pinger:]wird zusammen mit dem Hintergrundrauschen ausgef"uhrt und ist um den Linux eigenen Ping Befehl %
aufgebaut. Mittels Pinger wird die Latenz unter den einzelnen Endger"aten gemessen. %
\end{description}
\end{enumerate}
\section{Einsatz im Netzwerk} 

Mit dem Einsatz der im vorherigen Abschnitt vorgestellten Software ist das Testframework aufgebaut. In der \cref{fig:AufbauVomNetzwerk} %
wird dargestellt, wie die einzelnen Komponenten zusammenspielen. Hier sieht man, dass an einem Switch zwei Raspberry Pis %
und an einem anderen Switch drei Raspberry Pis angeschlossen sind. Einer von diesen Pis ist der Zabbix Server, der %
die Aktiven Hosts im Netzwerk "uberwacht. 

\label{sec:einsatzImNetzwerk}
\begin{figure}[htbp]
\centering
\includegraphics*[width=0.9\linewidth]{Abb/Netzschaltung3}

\caption{Aufbau des Netzwerks}
\label{fig:AufbauVomNetzwerk}
\end{figure}















%----------------- KAPITEL : TESTAUFBAU  ----------------- %
\chapter{Zabbix}
\label{cha:zabbix}
In \cref{cha:framework} wurde schon die verwendete Netzwerk Monitoring Software angeschnitten. In diesem Abschnitt wird nun %
ein tieferer Einblick in den Aufbau und die Verwendung von Zabbix gew"ahrt. Zabbix erm"oglicht es %
auch das Monitoring als ein Verteiltes Monitoring aufzuziehen mit mehreren Servern, diese sogenannten Proxies %
sind jedoch kein Bestandteil dieser Thesis und werden deshalb nicht weiter betrachtet.

\section{Zabbix Server}
\label{sec:server}
Der Zabbix Server ist die Zentrale Komponente des vorgestellten Frameworks. Dem Server werden von den Agents und Proxies %
Informationen zugeschickt. Die Aufgabe des Server ist es die Daten zu speichern und zu verarbeiten. %
Der Server speichert auch die Konfiguration der einzelnen Agents die sich im Netzwerk befinden. %
Um die Konfiguration der Agents einstellen zu k"onnen m"ussen jedoch erst einmal die Agents im Zabbix %
Server registriert werden. Die dazu ben"otigten Daten sind die Agent IP und ein eindeutiger Name. %
Wenn nun der Agent im Server registriert worden ist kann man im Zabbix Server den Agents verschiedene Templates auflegen. %
Das in dem in \cref{cha:framework} gezeigte Framework verwendet prim"ar das Template: \emph{Template OS %
Linux}. Dieses Templates ist eines von vielen vordefinierten Templates. Es ist auch m"oglich selber Templates %
zu erstellen, Templates sind in einem XML Format abgespeichert und k"onnen dadurch einfach unter Nutzern von Zabbix % 
ausgetauscht werden. Die Firma die den vertrieb von Zabbix regelt hat extra daf"ur die Zabbix Share eingef"uhrt %
auf dem Zabbix Templates und vieles mehr ausgetauscht werden kann \autocite{zabbix:share}. \ \\ %
Templates enthalten sogenannte Items die "uber Active Checks oder Passive Checks Informationen von %
einem vorher im Zabbix Server registrierten Agent anfordern. Ein Passive Check geschieht "uber einen %
Request diese sind in JSON verfasst und werden in geb"undelt verschickt um Bandbreite zu sparen. %
Der Prozess beeinhaltet f"unf Schritte: %
\begin{enumerate}
\item Der Server "offnet eine TCP Verbindung %
\item Der Server sendet einen request als Beispiel agent.version %
\item Der Agent empf"angt den request und antwortet mit der Versionsnummer %
\item Der Server verarbeitet die Antwort und erh"alt als Ergebniss Zabbix3.0.0rc %
\item Die TCP Verbindung wird geschlossen %
\end{enumerate}
\autocite{zabbix:activepassive}. Passive Checks und Active Checks unterscheiden sich darin, wer den request %
ausl"ost. Bei einem passiv Check sind die Hosts selber inaktiv bis vom Server ein request eintrifft, bei Active %
Checks wird der Agent selber Aktiv. Der Agent selber sendet dann an den Server einen Request in dem der Agent nach den %
Daten fragt die der Server erwartet sind\autocite{zabbix:activepassive2}. %
Dies ist der Fall wenn vom Agent Daten abgefragt werden sollen die nicht vom %
Betriebssystem des Agents gesammelt werden. Das bedeutet das der Agent selber nun die Daten die der Server erwartet %
sammeln muss, dies geschieht meistens "uber externe Programme, sollte ein Agent zu viele Active Checks haben kann es %
sich auf die Performanz des Hosts auswirken. \ \\ % 

"Uber die API ist es m"oglich Programme f"ur den Server zu schreiben mit der man zum Beispiel %
eine eigene Benutzeroberfl"ache erstellen kann um die Konfigurationen auf dem Server zu verwalten. %
"Uber die Api ist es auch m"oglich einen eigenen Zugriff auf die dem Server zugrunde liegende %
Datenbank herzustellen. %
Der Zabbix Server basiert auf einem Apache Web Server. Um das Zabbix Frontend nutzen zu k"onnen wird %
PHP 5.4.0 ben"otigt, jedoch funktioniert PHP 7 noch nicht. Die Daten und Statistiken die Zabbix sammelt %
werden in einer Datenbank gespeichert. Es werden f"unf verschiedene Datenbanken unterst"utzt. %
\begin{itemize}
\item MySql Version 5.0.3 aufw"arts
\item Oracle Version 10g aufw"arts
\item PostgreSQL Version 8.1 aufw"arts
\item SQLite Version 3.3.5 aufw"arts
\item IBM DB2 Version 9.7 aufw"arts (Noch nicht fehlerfrei)
\end{itemize}
\autocite{zabbix:req}
Der Zabbix Server selber ist auch als ein Agent konfiguriert, der sich selber "uberwacht. Jedoch ist der Server %
nicht im allgemeinen Prozess des in \cref{cha:framework} vorgestellten Frameworks t"atig. %
Trotzdem wird in den folgenden Kapiteln der Zabbix Server nicht aussen vor gelassen und in die %
den Tests betrachtet werden, dabei werden haupts"achlich die von Zabbix Statistiken zur Perfomance und zum %
Traffic auf dem Ethernet Port des Servers genommen. Der Wert Perfomance betrachtet zwei Dinge, einmal die sogenannte Queue, welche %
angibt wie viele von den Agents "ubermittelten Werten darauf warten vom Zabbix Server verarbeitet zu werden und die %
die verarbeiteten Werte pro Sekunde. Diese Werte lassen einen Schluss auf die Leistungsanforderungen des Servers schlie"sen. %
Sollte die Queue einen bestimmten schwellwert erreichen wird auf dem Zabbix Frontend eine Warnung ausgegeben das der %
Server mit der Verarbeitung nicht hinterherkommt. Dies kann soweit gehen das die gesammelten Daten auf dem Server fehlerhaft %
sind. Jedoch ist dieses Szenario in diesem Framework unwahrscheinlich, da die gr"o"se des %
Frameworks "uberschaubar ist. Der Traffic auf dem Ethernet Port Eth0 wird betrachtet, da alle von den Agents gesammelten Informationen %
"uber diesen an den Server gelangen. Dabei wird der eingehende sowie auch der ausgehende Traffic betrachtet.

\section{Zabbix Agent}
\label{sec:agent}
Der Zabbix Agent wird auf dem zu "uberwachenden System installiert. Der Agent sammelt die Daten "uber im Betriebssystem %
integrierte Monitoring Funktionen oder "uber die Zabbix eigene API und sendet diese je nach Art des Checks diese Informationen an den Agent. %
Da der Agent schon im Betriebssystem integrierte Funktionen benutzt ist dieser sehr effizient und verbraucht kaum sp"urbar %
die Ressourcen des Host Systems. Anders als der Server besitzt der Agent kein eigenes Frontend und speichert die Werte nicht %
in einer Datenbank. Der Agent ist so konstruiert das dieser weitesgehend ohne Administrator Rechte funktionieren kann. % 
Da es auch m"oglich ist "uber einen Fernzugriff Agent Hostsysteme auszuschalten oder neuzustarten m"ussen f"ur diese Funktionen %
dem Agent in der Rechteverwaltung diese Rechte zugeteilt werden. Dadurch wird das Sicherheits%
risiko f"ur den Host erheblich reduziert. F"ur gew"ohnlich wird f"ur den Agent ein eigener User angelegt. Die Anwendung des Agents %
l"auft unter Unix Systemen als ein Daemon und unter Windows als ein Systemdienst. Zabbix unterst"utzt jedoch noch mehr %
Betriebssysteme eine Auswahl davon w"aren: 
\begin{itemize}
\item Linux
\item Windows: alle Desktop und Server Versionen seit 2000
\item OpenBSD
\item Mac OS X
\item Solaris 9,10,11
\end{itemize}
\autocite{zabbix:agent}.
Im Aufbau des Frameworks das in dieser Arbeit weiter betrachtet wird laufen die Zabbix Agents unter dem g"angigen Raspberry Pi %
Betriebssystem Raspbian, welches ein Debian Port ist. Um den Host jedoch verwenden zu k"onnen muss man in %
den Konfigurations Dateien die Ports die der Server f"ur das versenden von Requests benutzt eingestellt werden und auch die IP des Servers %
muss eingetragen werden. Da der Agent sonst eingehende Requests verwirft und somit die Sicherheit f"ur den Host gew"ahrleistet. %
In den Konfigurationsdateien der Agents ist es auch m"oglich %
die Active Checks zu notieren, was in diesem Versuchsaufbau auch gemacht wurde. Da das Template OS Linux keine Items hatte die das %
"uberwachen von der Festplatten Last hatte kann man dies "uber die Konfigurationsdateien einstellen. Das sieht in der Konfigurationsdatei %
dann so aus.
\begin{quote}

UserParameter=custom.vfs.dev.read.ops[*],customParaScript.sh \$1 4 

\end{quote}

Wenn dann in der Antwort zu einem Active Check nach dem \emph{custom.vfs.dev.read.ops[*]} gefragt wird, wird das Bash Script customParaScript.sh %
ausgef"uhrt. Die Auszeichnung emph{UserParameter=} gibt an das es sich hierbei um ein selbstdefiniertes Item handelt, welches mithilfe von Zabbix %
"uberwacht werden soll. Das Script customParaScript.sh dient dazu die Last auf der Festplatte nachzupr"ufen %
dazu liest es aus der Datei die in Unix Systemen unter /proc/diskstats liegt, nach den zugeh"origen Werten \$1 und 4 sind Variablen die f"ur die %
verarbeitung notwendig sind. Der in diesem Beispiel genannte UserParameter liest die Lese Operationen pro Sekunde aus der Datei und leitet diese %
an den Server weiter.


%----------------- KAPITEL : VERSUCHE  ----------------- %
\chapter{Versuche}
\label{cha:versuche}
Die beste Art ein Netzwerktestframework zu Testen ist in dem man das Framework benutzt %
dazu habe ich am Anfang der Arbeiten selber mehrere Tests definiert, jeder dieser Tests %
ist eine in einem Netzwerk m"ogliche Fehlerquelle die die Leistung des Netzwerkes %
beeintr"achtigen kann. Um Fehler identifizieren %
zu k"onnen m"ussen jedoch auch die Normalwerte bestimmt werden. Dazu muss erstmal ein %
St"orungsfreies Netzwerk aufgebaut werden und dieses muss beobachtet werden. %
Die in diesem Netz gesammelten Werte stellen die Basis f"ur unsere sp"ateren Ergebnisse dar. %
Der zweite Schritt ist es das Netz mit Fehlern aufzubauen und zu beobachten. Dadurch %
haben wir die Werte eines Fehlerbehafteten Netzwerkes und k"onnen diese nun mit dem %
fehlerfreien Netzwerk vergleichen. Dadurch sollte es uns nun m"oglich sein automatisierte %
Aussagen "uber die Art des Fehlers in einem Netzwerk zu treffen. %

\section{Versuchsbeschreibung}
\label{sec:versuchbeschreibung}
\begin{description}
\item[Ethernetkabel ohne Isolierung] Ethernetkabel werden oft sehr strapaziert %
dies kann dazu fuhren das sich die Isolation l"ost und somit Strahlung einwirkt. %
Dies kann die Leistung im Netzwerk beeintr"achtigen. 
\item[Falsch gedrehtes Twisted Pair Kabel] TP Kabel sind die am meisten in einem Ethnernet %
Netzwerk vorkommenden Kabel, es entsteht eine hohe St"orungsicherheit bei TP Kabeln %
sollten jedoch TP Kabel entdrillt worden sein kann eine hohe Strahlenbelastung auftreten. %
\item[Loop] Ein Loop ist wenn ein Switch mit sich selbst verbunden ist, der Switch sendet %
dann st"andig Datagramme an sich selbst und blockiert so jeden Traffic auf der Leitung. %
\item[Nicht angeschlossenes Kabel] Wie verh"alt sich das Netzwerk wenn ein Kabel von einem %
Endger"at im laufenden Betrieb entkoppelt wird.
\item[Forkbomb] Die Forkbomb ist ein Programm das von sich rekursiv Kopien erstellt %
so das der Computer alle sein Resourcen verbraucht nur noch diesen Code auszufuhren. %
\begin{quote} 
:()\{ :|:\&\};: [https://de.wikipedia.org/wiki/Forkbomb] 
\end{quote}
\item[Festplatte l"auft voll] In einem Netzwerk werden st"andig Daten versand, was also %
passiert wenn die Festplatte von einem Rechner voll"auft und der PC somit Arbeitsunf"ahig wird. %
\item[IP Adresse doppelt belegt] Firmennetze benutzen statische IPs damit Webserver immer unter %
dem selben Namen erreichbar bleiben, das jedoch zu konfigurieren kann zu Fehlern Menschlichen %
Fehlern fuhren, was passiert in einem Netz wenn zwei Rechner sich eine IP teilen. %
\item[Kollisionsdom"anen] Kollisionsdom"anen sind bereiche in einem Netzwerk wo sich mehrere Rechner eine %
eine Leitung teilen, dies ist geschieht in der Physikalischen Ebene eines Netwerks. Die Frage ist kann
man eine Kollisionsdom"ane erkennen. 
\end{description}
\section{Raspberry Pi Versuche}
\label{sec:raspberryPiVersuche}
\subsection{Gemachte Tests}
\subsection{Ergebnisse}

\section{Virtual Machine Versuche}
\label{sec:VMVersuche}
\subsection{Gemachte Tests}
\subsection{Ergebnisse}



%----------------- KAPITEL : VERGLEICH------------ %	
\chapter{Vergleich VM/HW}
\label{cha:vergleich}
\section{Versuchsergebnisse}
\section{Kosten nutzen Faktor}


%----------------- KAPITEL : FAZIT  ----------------- %	
\chapter{Fazit}
%\chapter{Anhang}
\label{cha:anhang}

\section{Test Ruhezustand}

\begin{figure}[htbp]
\centering
\includegraphics*[width=1\linewidth]{Abb/ZabbixDazzle/NoTraffic/DazzleNoTrafficEth0}

\caption{Traffic auf Eth0 ohne Hintergrundrauschen auf Dazzle.}
\label{fig:Eth0DazzleNotraffic}
\end{figure}


\begin{figure}[htbp]
\centering
\includegraphics*[width=0.9\linewidth]{Abb/ZabbixDazzle/NoTraffic/DazzleNoTrafficIoStat}

\caption{I/O Statistik von Dazzle ohne Hintergrundrauschen.}
\label{fig:IoStatDazzleNotraffic}
\end{figure} %

\begin{figure}[htbp]
\centering
\includegraphics*[width=0.9\linewidth]{Abb/ZabbixDazzle/NoTraffic/DazzleNoTrafficCPU}

\caption{Prozentuale Lastverteilung auf der CPU von Dazzle, ohne Hintergrundrauschen.}
\label{fig:CPUDazzleNotraffic}
\end{figure}
\FloatBarrier
\section{Test 20 Megabyte Tests}

\begin{figure}[htbp]
\centering
\includegraphics*[width=0.9\linewidth]{Abb/ZabbixDazzle/Standard/Standard20}
\caption{Traffic auf Eth0 bei 20 Megabyte Dateien auf Dazzle.}
\label{fig:Standard20}
\end{figure}


\begin{figure}[htbp]
\centering
\includegraphics*[width=0.9\linewidth]{Abb/ZabbixDazzle/Standard/IoStatStandard20}

\caption{I/O Statistik von Dazzle beim 20 Megabyte Betrieb.}
\label{fig:IoStatStandard20}
\end{figure} %

\begin{figure}[htbp]
\centering
\includegraphics*[width=0.9\linewidth]{Abb/ZabbixDazzle/Standard/CPUStandard20}

\caption{Prozentuale Lastverteilung auf der CPU bei 20 Megabyte Paketen von Dazzle.}
\label{fig:CPUDazzleStandard}
\end{figure}
\FloatBarrier
\newpage
\section{200 Megabyte Tests}

\begin{figure}[htbp]
\centering
\includegraphics*[width=0.9\linewidth]{Abb/ZabbixDazzle/Standard/Standard200}
\caption{Traffic auf Eth0 bei 200 Megabyte Dateien auf Dazzle.}
\label{fig:Standard200}
\end{figure} %


\begin{figure}[htbp]
\centering
\includegraphics*[width=0.9\linewidth]{Abb/ZabbixDazzle/Standard/IoStatStandard200}

\caption{I/O Statistik von Dazzle beim 200 Megabyte Betrieb.}
\label{fig:IoStatDazzleStandard200}
\end{figure} %

\begin{figure}[htbp]
\centering
\includegraphics*[width=0.9\linewidth]{Abb/ZabbixDazzle/Standard/CPUStandard200}

\caption{Prozentuale Lastverteilung auf der CPU bei 200 Megabyte Paketen von Dazzle.}
\label{fig:CPUDazzleStandard200}
\end{figure}
\FloatBarrier
\newpage
\section{2000 Megabyte Pakete}

\begin{figure}[htbp]
\centering
\includegraphics*[width=0.9\linewidth]{Abb/ZabbixTusk/Standard/Standard2000}
\caption{Traffic auf Eth0 bei 2000 Megabyte Dateien auf Dazzle.}
\label{fig:Standard2000}
\end{figure} %


\begin{figure}[htbp]
\centering
\includegraphics*[width=0.9\linewidth]{Abb/ZabbixDazzle/Standard/IoStatStandard2000}

\caption{I/O Statistik von Dazzle beim 2000 Megabyte Betrieb.}
\label{fig:IoStatDazzleStandard2000}
\end{figure} %

\begin{figure}[htbp]
\centering
\includegraphics*[width=0.9\linewidth]{Abb/ZabbixDazzle/Standard/CPUStandard2000}

\caption{Prozentuale Lastverteilung auf der CPU bei 2000 Megabyte Paketen von Dazzle.}
\label{fig:CPUDazzleStandard2000}
\end{figure}

\section{Inhalt der CD}
Auf der CD sind folgende Inhalte:
\begin{itemize}
\item Die Logfiles zu den Tests, 
\item die selbsterstellten Programme,
\item die Graphen der restlichen Hosts,
\item eine Libre Office Calc Tabelle.
\end{itemize}




















%----------------- KAPITEL : LITERATUR  ----------------- %	
\printbibliography
%----------------- KAPITEL : ANHANG  ----------------- %	


\end{document}
